% LANA_REVISION_DRAFT_20260708T140659Z
% Paper: M2 P1 / m2_p1_outflow_escape_recycling
% This is a lane-local revision draft only. It does not overwrite the current linked manuscript or PDF.
\documentclass[twocolumn]{aastex631}

\usepackage{amsmath}
\usepackage{booktabs}
\graphicspath{{/Users/duhokim/NebulaMind/NebulaMind/.hermes/handoffs/galaxy-evolution/mastermind/aas-autopilot/runs/SDSS_REMAINING_TOPIC_PILOTS_20260708T125828Z/m2_p1_outflow_escape_recycling/figures/}}

\shorttitle{High-Excitation Optical AGN Denominator}
\shortauthors{NebulaMind Autopilot}

\begin{document}

\title{A High-Excitation Optical-AGN Denominator for Escape-versus-Recycling Follow-up in SDSS DR17}

\author{NebulaMind Research Autopilot}
\affiliation{Local reproducible pilot run; public SDSS DR17 data only}

\begin{abstract}
This revision narrows the outflow escape/recycling pilot to what the SDSS-only analysis actually measures: a high-excitation optical-AGN denominator for future resolved-kinematic follow-up.  In the cached 60,000-galaxy SDSS DR17 emission-line sample, the batch analysis identifies 4,440 high-excitation optical AGN candidates, a fraction of 0.074 with approximate binomial 95\% interval [0.0719, 0.0761].  Their median catalog $\log({\rm sSFR}/{\rm yr}^{-1})$ is -11.53, compared with -10.14 for the full denominator, a median offset of -1.39 dex.  The result is useful for selecting a feasible parent population for ionized, molecular, neutral, and CGM follow-up.  It is not a measurement of outflow velocity, escape velocity, mass loading, recycling time, or gas escape fraction.
\end{abstract}

\keywords{galaxies: active --- galaxies: evolution --- galaxies: outflows --- surveys --- methods: data analysis}

\section{Purpose and Scope}\label{sec:scope}
The full research topic asks whether AGN-driven multiphase outflows escape galaxy halos or recycle back into the circumgalactic medium.  The present SDSS pilot contains optical line ratios and catalog galaxy properties, but it does not contain spatially resolved outflow velocities, halo escape velocities, molecular or neutral gas masses, or CGM absorption/recycling diagnostics.  The manuscript should therefore be framed as a target-denominator paper.

The useful pilot question is: how large is the high-excitation optical-AGN subset in the public SDSS emission-line sample, and how does its catalog sSFR compare with the parent denominator?  That question is answerable from the preserved analysis summary and provides a bridge to the follow-up data required by the full proposal.

\section{Data and Candidate Definition}\label{sec:data}
The parent sample is the same 60,000-row SDSS DR17 emission-line denominator used across the overnight pilot set.  The selection requires spectroscopic galaxy classification, $0.02<z<0.12$, finite stellar-mass and sSFR estimates, and signal-to-noise ratio at least 3 in H$\alpha$, H$\beta$, [O~III] $\lambda5007$, and [N~II] $\lambda6584$.  BPT classes are recomputed from line ratios using the Kauffmann et al. and Kewley et al. demarcations.

The run summary reports a high-excitation optical-AGN candidate subset of 4,440 galaxies.  Before this draft is merged, the integration pass should insert the exact line-ratio or BPT-subclass criterion from the preserved analysis code.  This revision deliberately avoids inventing a threshold that is not present in the JSON summary.

\section{Results}\label{sec:results}
Table~\ref{tab:m2p1_summary} converts the current itemized result into a manuscript-ready quantitative summary.  The candidate subset is large enough for follow-up planning, and it is substantially lower in median sSFR than the full emission-line denominator.  Neither fact demonstrates escape or recycling; both facts define a practical parent sample.

\begin{deluxetable*}{lcc}
\tablecaption{High-excitation optical-AGN denominator summary\label{tab:m2p1_summary}}
\tablehead{\colhead{Quantity} & \colhead{Value} & \colhead{Interpretation}}
\startdata
Parent SDSS emission-line denominator & 60,000 galaxies & Shared pilot sample \\
High-excitation optical AGN candidates & 4,440 galaxies & Follow-up target pool \\
Candidate fraction & 0.074 $\pm$ 0.001 & Approx. 95\% interval [0.0719, 0.0761] \\
Median $\log {\rm sSFR}$, full denominator & -10.14 & Catalog sSFR baseline \\
Median $\log {\rm sSFR}$, candidates & -11.53 & Lower-sSFR target subset \\
Median candidate offset from full denominator & -1.39 dex & Descriptive only; not causal feedback \\
\enddata
\tablecomments{Counts, fractions, and median sSFR values are copied from the preserved topic-specific JSON summary.  The candidate-definition row should be made fully explicit from the analysis code before integration.}
\end{deluxetable*}

\begin{figure}
\centering
\includegraphics[width=\columnwidth]{m2\_p1\_outflow\_escape\_recycling\_figure1.pdf}
\caption{Preserved SDSS DR17 pilot diagnostic for the high-excitation optical-AGN denominator.  The figure should be captioned as a target-selection baseline; SDSS single-fiber optical data do not determine whether gas escapes or recycles.}
\label{fig:m2p1_denominator}
\end{figure}

\section{Interpretation Guard}\label{sec:guard}
The draft should explicitly separate three layers of inference.
\begin{enumerate}
\item \textbf{Measured here:} an optical candidate count, candidate fraction, and catalog sSFR comparison within one SDSS denominator.
\item \textbf{Needed for escape:} resolved outflow velocities, geometry, inclination, gravitational potential or escape velocity, and phase-dependent mass-loading estimates.
\item \textbf{Needed for recycling:} CGM or halo-gas tracers, timescale modeling, and multiphase gas measurements that connect launched material to reaccretion or retention.
\end{enumerate}

A safe discussion sentence is: ``The SDSS pilot identifies a feasible high-excitation optical-AGN parent sample for escape/recycling follow-up, but it does not distinguish escaping from recycling outflows.''  Unsafe wording would say that the pilot measures escape, recycling, or feedback efficiency.

\section{Discussion Outline for Integration}\label{sec:discussion}
The integrated manuscript can use the pilot as a staging paper.  First, define the denominator and high-excitation candidate criterion transparently.  Second, state that the observed lower median sSFR motivates, but does not prove, a feedback connection.  Third, propose concrete follow-up observables: IFU ionized-gas kinematics, molecular outflow tracers, neutral absorption, halo-mass or rotation-curve estimates for escape speeds, and CGM absorption or emission for recycling.  Fourth, preserve nondetections because they constrain how common measurable outflows are in the candidate denominator.

\section{Conclusions}\label{sec:conclusions}
\begin{enumerate}
\item The revised M2 P1 manuscript should be a denominator paper for escape/recycling follow-up, not an outflow-fate result.
\item The cached SDSS pilot identifies 4,440 high-excitation optical AGN candidates among 60,000 emission-line galaxies, giving fraction 0.074.
\item The candidates have median catalog $\log {\rm sSFR}=-11.53$, 1.39 dex below the full-denominator median.
\item The exact high-excitation criterion must be inserted from the preserved analysis code before integration; the current JSON summary alone is insufficient for a final methods section.
\end{enumerate}

\section*{Reproducibility and Safety Note}
Source analysis summary: \texttt{m2\_p1\_outflow\_escape\_recycling/analysis\_results.json} under run SDSS\_REMAINING\_TOPIC\_PILOTS\_20260708T125828Z.  This Lana revision is lane-local and did not overwrite the public-linked manuscript or PDF.

\acknowledgments
This pilot used public SDSS DR17 data products and open-source Python tools.

\begin{thebibliography}{}
\bibitem[Baldwin et al.(1981)]{baldwin1981} Baldwin, J.~A., Phillips, M.~M., \& Terlevich, R. 1981, PASP, 93, 5
\bibitem[Kauffmann et al.(2003)]{kauffmann2003} Kauffmann, G., Heckman, T.~M., Tremonti, C., et al. 2003, MNRAS, 346, 1055
\bibitem[Kewley et al.(2001)]{kewley2001} Kewley, L.~J., Dopita, M.~A., Sutherland, R.~S., Heisler, C.~A., \& Trevena, J. 2001, ApJ, 556, 121
\bibitem[York et al.(2000)]{york2000} York, D.~G., Adelman, J., Anderson, J.~E., Jr., et al. 2000, AJ, 120, 1579
\end{thebibliography}

\end{document}
